% problem_id: S047-b21-Q3
% source_id: S047/S048 (msp/pjm-305-1-p01.pdf)
% source_file: pjm-305-1-p01.pdf
% statement_locator: printed p. PDF p. 24 = printed p. 23. The statement is the last item on the page and is introduced by the author's lead-in sentence "The final lemma needed to establish the theorem can be proved in a similar manner to the lemma above."
% proof_locator: printed p. PDF pp. 24-25 = printed pp. 23-24. The proof begins on PDF p. 24 immediately after the statement and runs to "Case 2b: Sigma has two components. This case is nearly identical to the previous." at the foot of PDF p. 24; it resumes at the top of PDF p. 25 with the two-component filling argument and ends at the end-of-proof square immediately above "Proof of Theorem A.1". Context read verbatim for the chain: PDF p. 17 (Theorem A.1), PDF p. 18 (definitions of Sym, Sym+, Z(X), primitive curve, Delta, orbi-solid torus, orbi-lens space, strong involution), PDF pp. 19-20 (Lemma A.3 statement), PDF p. 22 (Lemma A.4 statement), PDF p. 23 (Lemma A.5, Lemma A.6 statements), PDF p. 25 (proof of Theorem A.1).
% representation: PDF; transcribed from rendered page images
% collected_by: carrier rule
% review_status: H2 by 2/2 independent reviewers (the miner claimed H3); 2/2 报不忠实
% status: complete (statement + author proof verbatim + reason + locators)
%
% =====================================================================
% Candidate: Q3
% Paper: Kenneth L. Baker (Appendix written jointly with Neil R. Hoffman),
%        "The Poincare homology sphere, lens space surgeries, and some knots
%        with tunnel number two", Pacific J. Math. 305 (2020), no. 1, 1-30.
%        DOI 10.2140/pjm.2020.305.1
% Source file: /disks/sata1/yupeng/human-proof-corpus/source-cache/registry-sources/msp/pjm-305-1-p01.pdf
%        (30 PDF pages; PDF page k == printed page k-1, i.e. PDF p. 24 = printed p. 23)
% Rendered with: pdftoppm -f <p> -l <p> -r 150 -png (whole pages, read as images)
%        plus targeted -r 300 / -r 600 crops for every symbol-bearing line of the unit
%        (bars, primes, Sigma, script Z, the subscript sigma'_{n,m}).
%        pdftotext -layout was used ONLY to locate items; no symbol was read from the text layer.
%
% ALREADY COLLECTED FROM THIS PAPER (deliberately NOT repeated here):
%        Theorem 1 (printed pp. 1-9) and Lemma 12 (printed p. 9).
%
% COLLECTED UNIT:  Lemma A.7 of the Appendix (Baker-Hoffman),
%        printed pp. 23-24 = PDF pp. 24-25, WITH its complete author proof.
%
% Context transcribed below, verbatim and clearly marked as context:
%        A. Theorem A.1 + its proof (the appendix flagship that Lemma A.7 carries);
%        B. the appendix's definitions of Sym(X), Sym+(X), Z(X), primitive curve,
%           distance Delta, orbi-solid torus, orbi-lens space, strong involution;
%        C. the statements of Lemma A.3, Lemma A.4, Lemma A.5, Lemma A.6
%           (the in-paper results cited by Lemma A.7's proof; their proofs are NOT
%           transcribed -- they are located by page in Q3.json).
%
% Verbatim transcription. No rewriting, no completion, no smoothing.
% Original typos are kept and flagged with [sic: ...]. Illegible spots are flagged
% as %% [ILLEGIBLE: ...].  Nothing below is flagged illegible: every line of the unit
% was read from a 300/600-dpi crop.
% =====================================================================

\documentclass[11pt]{article}
\usepackage{amsmath}
\usepackage{amssymb}
\usepackage{amsthm}
\usepackage[margin=1in]{geometry}
\newcommand{\Sym}{\ensuremath{\mathrm{Sym}}}
\newcommand{\Zs}{\ensuremath{\mathcal{Z}}}   % the quotient orbifold Z = X/Z(X)
\newcommand{\Ps}{\ensuremath{\mathcal{P}}}   % the Poincare homology sphere
\newcommand{\Ns}{\ensuremath{\mathcal{N}}}   % regular neighbourhood
\begin{document}

\section*{Q3 --- Lemma A.7 of the Appendix (verbatim transcription)}

\noindent\textbf{Source.} K.~L.~Baker, with an appendix written jointly with N.~R.~Hoffman,
\emph{The Poincar\'e homology sphere, lens space surgeries, and some knots with tunnel number two},
Pacific J.\ Math.\ \textbf{305} (2020), no.~1, 1--30. Appendix, printed pp.~23--24.

% ---------------------------------------------------------------------
% LEAD-IN SENTENCE (printed p. 23 = PDF p. 24, verbatim): the last sentence of
% the preceding paragraph, which introduces the lemma.  It is transcribed because
% it is the author's own statement about how the lemma relates to Lemma A.6.
% ---------------------------------------------------------------------
The final lemma needed to establish the theorem can be proved in a similar
manner to the lemma above.

%% --- page 23 (PDF p. 24) ---

\subsection*{Lemma A.7 (printed p.~23 = PDF p.~24)}

\begin{quote}
\textbf{Lemma A.7.} \emph{Let $X$ be the exterior of a hyperbolic knot $K$ in a lens space
$L(p,q)$ with $p \ge 2$ admitting a Dehn surgery to $\Ps$. Then $|Z(X)| = 1$.}
\end{quote}

\subsection*{Proof (authors' text, verbatim)}

\textit{Proof.} Set $Z = Z(X)$ and assume $|Z| > 1$. Let $\mu$ be the meridian of $K$ and let
$\gamma$ be the slope of the surgery to $\Ps$. Then let $\bar\mu$ and $\bar\gamma$ be the images
of $\mu$ and $\gamma$ in the quotient $\Zs = X/Z$.

By Lemma A.6, $\Zs$ is an orbifold (with nonempty singular set). Since $\Ps$ is a homology
sphere, $\Zs$ is the exterior of a knot $\bar K$ in an orbi-lens space $\Zs(\bar\mu)$ that is
disjoint from the singular set by Lemma A.3. Because $\Zs$ is a suborbifold of $\Zs(\bar\mu)$,
the singular set $\Sigma(\Zs)$ may be identified with the singular set $\Sigma(\Zs(\bar\mu))$.
Thus $\Sigma(\Zs)$ has at most two components. Since $\Ps = X(\gamma)$, the quotient orbifold
$\Ps/Z = \Zs(\bar\gamma)$ has a singular set $\Sigma(\Zs(\bar\gamma))$ consisting of one, two, or
three embedded circles according to whether the singular set $\Sigma(\Zs)$ has one or two
components and whether or not the slope $\bar\gamma$ is a primitive curve.

\textit{Case 1: $\bar\gamma$ is not primitive.} Then $\Zs$ is the exterior of one of the
components of $\Sigma(\Zs(\bar\gamma))$. Hence by Lemmas A.4 and A.5, $\Zs$ must be Seifert
fibered. Thus $X$ must be Seifert fibered contrary to $K$ being a hyperbolic knot.

\textit{Case 2: $\bar\gamma$ is primitive.} Let $\bar K_\gamma$ be the core of the filling solid
torus in $\Zs(\bar\gamma)$ so that the exterior of $\bar K_\gamma$ is $\Zs$. Since $\bar\gamma$
is primitive, the singular set $\Sigma = \Sigma(\Zs(\bar\gamma))$ of $\Zs(\bar\gamma)$ is contained
in the knot exterior $\Zs$ and thus has either one or two components. Since $\Zs$ is a hyperbolic
orbifold in which $\Sigma$ is a geodesic link of one or two components, the exterior of $\Sigma$
is a hyperbolic manifold $X' = \Zs_\Sigma = \Zs - \Ns(\Sigma)$ [Kojima 1988; Sakai 1991].

Let us write $X'(\alpha,\beta)$ (or $X'(\alpha,\beta_1,\beta_2)$) to denote the Dehn filling of
$X'$ along the $\alpha$ and $\beta$ (or $\alpha$ and $\beta_1$, $\beta_2$) slopes in the components
of $\partial X'$ corresponding to $\partial\Ns(\bar K_\gamma) = \partial\Zs$ and
$\partial\Ns(\Sigma)$ respectively. We will use $-$ to indicate that the boundary component is
left unfilled.

\textit{Case 2a: $\Sigma$ has just one component.} By Lemma A.4, $X'(\gamma,-)$
%% [sic: the printed text drops the bar over gamma here; the very next sentence prints
%%  X'(gamma-bar, -).  Transcribed as printed.]
is a Seifert fibered space over the disk with two exceptional fibers. Thus $X'(\bar\gamma,-)$ has
a one-parameter family of slopes $\sigma_i$ giving lens space fillings $X'(\bar\gamma,\sigma_i)$.
Since $X'(\bar\mu,-)$ is a solid torus (as it is the complement of the core of an orbi-solid torus
in an orbifold lens space with connected singular set), $X'(\bar\mu,\beta_i)$ is also a lens space
for each $\beta_i$. Hence, by Thurston's hyperbolic dehn filling theorem [1979, Theorem 5.8.2],
we can choose $\beta_i$ so that $X'(-,\beta_i)$ is hyperbolic. However,
$\Delta(\bar\mu,\bar\gamma) = |Z|\Delta(\mu,\gamma) \ge 2$, which contradicts the cyclic surgery
theorem [Culler et al. 1987].

\textit{Case 2b: $\Sigma$ has two components.} This case is nearly identical to the previous.

%% --- page 24 (PDF p. 25) ---

By Lemma A.5, $X'(\bar\gamma,-,-)$ is a Seifert fibered space over the annulus with one exceptional
fiber. Then there are infinitely many pairs of slopes
$\{(\sigma_n,\sigma'_{n,m})\}_{(n,m)\in\mathbb{Z}^2}$ such that
$X'(\bar\gamma,\sigma_n,\sigma'_{n,m})$ is a lens space. Since $X'(\bar\mu,-,-)$ is a thickened
torus (being the complement of the two cores of the orbi-solid tori in the orbi-lens space),
$X(\bar\mu,\sigma_n,\sigma'_{n,m})$
%% [sic: the printed text drops the prime on X here; the paragraph's notation is X'(...).
%%  Transcribed as printed.]
is a lens space for all these pairs of slopes.

Using Thurston's hyperbolic Dehn filling theorem twice, there are infinitely many of these pairs
of slopes such that $X(-,\sigma_n,\sigma'_{n,m})$
%% [sic: prime dropped on X again, as printed.]
is hyperbolic. Again, since $\Delta(\bar\mu,\bar\gamma) = |Z|\Delta(\mu,\gamma) \ge 2$, we obtain a
contradiction to the cyclic surgery theorem [Culler et al. 1987] for these pairs.
\hfill$\square$

% ---------------------------------------------------------------------
% TRAILING PARAGRAPH (printed p. 24 = PDF p. 25, verbatim).  It stands AFTER the
% end-of-proof square, so it is a remark, not part of the proof.  Transcribed
% because it states the argument's own limits.
% ---------------------------------------------------------------------
\medskip
\noindent\textit{[Trailing remark, printed immediately after the end-of-proof mark; NOT part of
the proof.]} It may be tempting to try to use $|Z|\cdot\Delta(\mu,\gamma)$ in the previous lemma
place try and improve the bound in [Boyer and Zhang 1996] that $\Delta(\mu,\gamma) \le 2$. However,
this bound is known to be sharp. Also, when $Z(M)$
%% [sic: printed Z(M); the appendix's notation for this group is Z(X).]
is trivial, we lose the ability to ``match up'' the cores of the solid tori coming from the Heegaard
splitting of the lens space with the exceptional fibers in $\Ps$, and so the arguments in this
appendix will not apply.

\bigskip
\hrule
\bigskip

\section*{Context A --- the flagship that Lemma A.7 carries (verbatim)}

% ---------------------------------------------------------------------
% Printed p. 16 = PDF p. 17 (statement) and printed p. 24 = PDF p. 25 (proof).
% ---------------------------------------------------------------------

\noindent\textbf{Theorem A.1.} \emph{Let $X$ be a one-cusped hyperbolic manifold admitting both a
Poincar\'e homology sphere filling and a lens space filling. Then the symmetry group of $X$ is
trivial or generated by a single strong involution.}

\medskip
\noindent\textit{Proof of Theorem A.1} (printed p.~24). By Proposition A.2,
\[
\Sym(X) = \Sym^+(X).
\]
From Lemma A.3(2), $Z(X)$ has index at most 2 in $\Sym(X)$. Then since $Z(X) = 1$ by Lemma A.7, we
see that $\Sym(X)$ is either trivial or $\mathbb{Z}/2\mathbb{Z}$. According to Lemma A.3(3), the
nontrivial element in this latter case must be a strong involution. \hfill$\square$

\medskip
\noindent\textit{[Appendix programme, printed pp.~16--17, verbatim; it locates Lemma A.7 as step
(c).]} Thereafter, the argument in this appendix will gradually ``whittle-down'' the symmetry group
of the exterior $X$ of a hyperbolic knot $K$ in a lens space with a Dehn surgery to $\Ps$. First,
we give an argument which eliminates orientation reversing symmetries so that
$\Sym(X) = \Sym^+(X)$. Then we proceed to analyze a subgroup $Z(X)$ of $\Sym^+(X)$, showing

(a) that it has index at most 2,

(b) if the index is 2 then $\Sym(X)$ contains a strong involution, and ultimately

(c) that $Z(X)$ is trivial.

This is pulled together in the proof of Theorem A.1 at the end of this appendix.

\bigskip
\hrule
\bigskip

\section*{Context B --- definitions used by the statement (verbatim)}

% ---------------------------------------------------------------------
% Printed pp. 17-18 = PDF pp. 18-19.  Omissions inside a quotation are marked
% explicitly with %% [omitted: ...]; nothing is reworded.
% ---------------------------------------------------------------------

\noindent\textit{[Printed p.~17, verbatim.]} \dots\ A \emph{primitive curve}
%% [omitted: the preceding definition of the singular set Sigma(Q), of the complement
%%  Q - K and the exterior X = Q - N(K), and the standing assumption that K is either
%%  a component of Sigma(Q) or disjoint from it.]
is homotopic to an essential simple closed curve, and hence is a slope. Any nonprimitive essential
curve is a multiple of a primitive curve. Given two closed curves $\alpha$ and $\beta$ in the torus
we say the \emph{distance} between $\alpha$ and $\beta$, $\Delta(\alpha,\beta)$ is the minimal
(unoriented) geometric intersection number between the two curves. If $\alpha$ is a curve in
$\partial\Ns(K)$ that is an $n$-fold multiple of a slope $\gamma$, then the result of $\alpha$-Dehn
surgery on $K \subset Q$ is the orbifold $X(\alpha)$ with underlying manifold $|X(\alpha)|$
obtained as the $\gamma$-Dehn surgery on $K \subset |Q|$ and a new singular set $\Sigma(X(\alpha))$
consisting of $\Sigma(X)$ and the core of the attached $S^1\times D^2$ with order $n$ (unless
$n = 1$ in which case the singular set remains only $\Sigma(X)$). The orbifold attached to $X$ is an
\emph{orbi-solid torus} of order $n$. An orbi-solid torus of order 1 is just a solid torus.

Following [Boileau et al.\ 2012, \S3], an \emph{orbi-lens space} is the quotient of $S^3$ by a
finite cyclic subgroup of $SO(4)$ and is consequentially a union of two orbi-solid tori with coprime
orders along their common boundary. Its singular set is the union of the cores of these orbi-solid
tori that have order at least 2.

\dots\ For a hyperbolic knot $K$ with exterior $X$, the symmetry group $\Sym(X)$ of homeomorphisms
modulo isotopy is identified with the isometry group of the interior of $X$. The orientation
preserving symmetry group is denoted $\Sym^+(X)$.

Within $\Sym^+(X)$ is the subgroup $Z(X)$, the maximal subgroup that acts freely on $\partial X$.
In particular, $Z(X)$ is the maximal subgroup of $\Sym^+(X)$ for which the quotient orbifold
(or manifold) $X/Z(X)$ has a torus cusp.

\dots\ Elements of $\Sym^+(X)$ of order 2 that reverse the orientation of the homological longitude
$\lambda$ of the manifold $X$ are called \emph{strong involutions}.

\bigskip
\hrule
\bigskip

\section*{Context C --- statements cited inside the proof of Lemma A.7 (verbatim)}

% ---------------------------------------------------------------------
% Statements only.  Their proofs live on printed pp. 18-20 (A.3), p. 21 (A.4) and
% p. 22 (A.4 proof end, A.5, A.6) = PDF pp. 19-21, 22, 23; see Q3.json.
% ---------------------------------------------------------------------

\noindent\textbf{Lemma A.3} (printed pp.~18--19 = PDF pp.~19--20). \emph{Let $X$ be the exterior of
a hyperbolic knot $K$ in a lens space $L(p,q)$ with $p \ge 2$ admitting a Dehn surgery to an
integral homology sphere. Then the following hold:}

(1) \emph{$Z(X)$ is cyclic.}

(2) \emph{$[\Sym^+(X) : Z(X)] \le 2$.}

(3) \emph{If $[\Sym^+(X) : Z(X)] = 2$, then $X$ has a symmetry which is a strong involution.}

(4) \emph{In the quotient orbifold $\Zs = X/Z(X)$, the meridian $\mu$ of $K$ maps to a primitive
curve $\bar\mu$ in $\partial\Zs$.}

(5) \emph{The Dehn filling $\Zs(\bar\mu)$ is an orbi-lens space in which the core $\bar K$ of the
filling is disjoint from the singular set of $\Zs(\bar\mu)$.}

\medskip
\noindent\textbf{Lemma A.4} (printed p.~21 = PDF p.~22). \emph{If the finite cyclic quotient of
$\Ps$ by symmetries is an orbifold $Q$, then $Q$ is homeomorphic to one of the following:}

(1) \emph{a manifold $M$ with $\pi_1(M) \cong \pi_1(\Ps) \times \mathbb{Z}/n\mathbb{Z}$ for some
integer $n$ coprime to 30,}

(2) \emph{an orbifold that fibers over $S^2(2,3,5)$ where the singular set of $Q$ has 1, 2, or 3
components, or}

(3) \emph{an orbifold with base space $S^3$ and singular set a knot or link as pictured in Figure
7.}

\emph{Moreover, in the case that $Q$ is an orbifold, the components of its singular set are fixed
locally by finite cyclic groups of relatively prime orders.}

\medskip
\noindent\textbf{Lemma A.5} (printed p.~22 = PDF p.~23). \emph{In Cases 2 and 3 of Lemma A.4:}

\begin{itemize}
\item \emph{If $|\Sigma(Q)| = 1$, then the complement of $\Sigma(Q)$ is a Seifert fibered space over
the disk with two exceptional fibers.}
\item \emph{If $|\Sigma(Q)| = 2$, then the complement of $\Sigma(Q)$ is a Seifert fibered space over
the annulus with one exceptional fiber.}
\item \emph{If $|\Sigma(Q)| = 3$, then the complement of $\Sigma(Q)$ is $F \times S^1$, where $F$ is
a pair of pants.}
\end{itemize}

\medskip
\noindent\textbf{Lemma A.6} (printed p.~22 = PDF p.~23) --- the ``lemma above'' of the lead-in
sentence. \emph{Let $X$ be the exterior of a hyperbolic knot $K$ in a lens space $L(p,q)$ with
$p \ge 2$ admitting a Dehn surgery to $\Ps$. Then every nontrivial subgroup of $Z(X)$ acts
nonfreely on $X$.}

\end{document}
